Hilbert（希尔伯特）的纲领旨在表明，所有数学都可以在一个形式语言（如算术语言或集合论语言）的公理化理论中得到形式化。
他相信这样一个理论将是\textbf{完备的}。即，对每个语句~$!A$，要么 $\Th{T} \Proves !A$，要么 $\Th{T} \Proves \lnot !A$。
这样，$\Th{T}$ 就将解决每一个数学问题：它要么证明其为真，要么证明其为假。
如果 Hilbert 是对的，那么数学也将是\textbf{可判定的}。
这是因为任何可公理化的理论都是\textbf{可计算枚举的}，即存在一个可计算函数能列出它的所有定理。
我们可以通过逐一列出所有定理来检验一个语句~$!A$ 是否为定理，直到找到 $!A$（此时它是定理）或找到 $\lnot !A$（此时它不是定理）。
唉，Hilbert 错了。G\"odel（哥德尔）证明了，没有哪个可公理化、一致且“足够强”的理论是完备的。这就是\textbf{第一不完备性定理}。
“足够强”的要求指的是该理论能够表示所有可计算函数和关系。具体来说，极弱的理论 $\Th{Q}$ 就满足这个性质，任何至少和~$\Th{Q}$ 一样强的理论也都满足。
他还证明了——即\textbf{第二不完备性定理}——表达该理论一致性的语句本身在该理论中是不可判定的，即该理论既不能证明它，也不能证明它的否定。
因此，Hilbert 寻找全部数学的“有穷”一致性证明的进一步目标也不可能实现。因为任何有穷的一致性证明，大概都可以在涵盖全部数学的理论中被形式化。
最后，我们还确立了，能够表示所有可计算函数和关系的理论是\textbf{不可判定的}。
请注意，虽然可公理化和完备性共同蕴涵可判定性，但不完备性并不蕴涵不可判定性。
所以，这个结果表明 Hilbert 的第二个目标——即存在一个能判定 $\Th{T} \Proves !A$ 是否成立的过程——也无法实现，至少对于不弱于~$\Th{Q}$ 的理论而言是如此。